\documentclass[../../main.tex]{subfiles}
\begin{document}

\begin{center}
  {\fontsize{48pt}{58pt}\selectfont \textbf{1966}\normalsize\textbf{年度}}
\end{center}

\begin{table}[htb]
  \centering
  \caption{出題テーマと難易度}
  \begin{tabular}{|c|p{9cm}|c|}
    \hline
    \textbf{問題番号} & \textbf{テーマ} & \textbf{難易度} \\
    \hline
    第1問 & 端数を切り上げる運賃計算 & \\
    \hline
    第2問 & 一次変換で不変な直線 & \\
    \hline
    第3問 & 等比的に伸びる折れ線の点の座標 & \\
    \hline
    第4問 & 三角方程式の最大の根 & \\
    \hline
    第5問 & 新課程：らせんの道のり／旧課程：積分条件と係数の範囲 & \\
    \hline
    第6問 & 新課程：四桁の数の確率／旧課程：放物線の鏡映 & \\
    \hline
  \end{tabular}
\end{table}

\begin{table}[htb]
  \centering
  \caption{解答}
  \renewcommand{\arraystretch}{1.4}
  \begin{tabular}{|c|c|p{9cm}|}
    \hline
    \textbf{問題番号} & \textbf{小問} & \textbf{答え} \\
    \hline
    第1問 & -- & $a=3.1$，$b=1.6$ \\
    \hline
    第2問 & -- & $y=-x$，$y=\dfrac{x}{2}$ \\
    \hline
    第3問 & -- & $P_{3n}=\left(\dfrac{2(1-2^{3n})}{7},\dfrac{\sqrt{3}(1-2^{3n})}{7}\right)$ \\
    \hline
    第4問 & -- & 最も近い整数は$5$ \\
    \hline
    第5問 & 新課程 & 道のりは$\dfrac{\sqrt{5}}{2}(e^{4\pi}-1)$ \\
    \cline{2-3}
    & 旧課程 & $a+b\le-4$，または$a+b\ge-\dfrac{4}{3}$ \\
    \hline
    第6問 & 新課程 & 確率は$\dfrac{23}{216}$ \\
    \cline{2-3}
    & 旧課程 & $l_1:y=3$，$l_2:y=-x+6$ \\
    \hline
  \end{tabular}
\end{table}

\paragraph{試験形態}
6問．第5・6問は新旧課程別．試験時間は要確認．

\end{document}
